\documentclass[11pt,a4paper]{article}
% preamble.tex — estilo común de todos los apuntes Froth.
% Cada apunte hace % preamble.tex — estilo común de todos los apuntes Froth.
% Cada apunte hace % preamble.tex — estilo común de todos los apuntes Froth.
% Cada apunte hace \input{preamble} justo después de \documentclass.
\usepackage[margin=2.4cm]{geometry}
\usepackage{xcolor}
\definecolor{litblue}{RGB}{31,79,121}
\definecolor{litgreen}{RGB}{27,94,32}
\definecolor{litorange}{RGB}{230,81,0}
\usepackage{parskip}
\usepackage{titlesec}
\titleformat{\section}{\Large\bfseries\color{litblue}}{\thesection}{0.6em}{}
\titleformat{\subsection}{\large\bfseries\color{litblue}}{\thesubsection}{0.6em}{}
\usepackage{tcolorbox}
\usepackage{fancyhdr}
\pagestyle{fancy}
\fancyhf{}
\fancyhead[L]{\small\color{litblue}Froth — Apuntes de ML}
\fancyhead[R]{\small\thepage}
\renewcommand{\headrulewidth}{0.4pt}
\usepackage[colorlinks=true,linkcolor=litblue,urlcolor=litblue]{hyperref}

% Cabecera de cada apunte: \cabecera{numero}{titulo}
\newcommand{\cabecera}[2]{%
  \begin{tcolorbox}[colback=litblue,coltext=white,colframe=litblue,arc=2mm]
    {\Large\bfseries Apunte #1 — #2}\\[3pt]
    {\small Proyecto Froth · aprender Machine Learning construyendo}
  \end{tcolorbox}\vspace{0.8em}}

% Cajas temáticas
\newtcolorbox{concepto}[1]{colback=blue!4,colframe=litblue,title={Concepto: #1},arc=1.5mm}
\newtcolorbox{practica}{colback=green!5,colframe=litgreen,title={Pruébalo tú (teclado en mano)},arc=1.5mm}
\newtcolorbox{ojo}{colback=orange!8,colframe=litorange,title={Ojo / error típico},arc=1.5mm}
\newtcolorbox{resumen}{colback=gray!8,colframe=gray!60,title={En una frase},arc=1.5mm}

\newcommand{\codigo}[1]{\texttt{#1}}
 justo después de \documentclass.
\usepackage[margin=2.4cm]{geometry}
\usepackage{xcolor}
\definecolor{litblue}{RGB}{31,79,121}
\definecolor{litgreen}{RGB}{27,94,32}
\definecolor{litorange}{RGB}{230,81,0}
\usepackage{parskip}
\usepackage{titlesec}
\titleformat{\section}{\Large\bfseries\color{litblue}}{\thesection}{0.6em}{}
\titleformat{\subsection}{\large\bfseries\color{litblue}}{\thesubsection}{0.6em}{}
\usepackage{tcolorbox}
\usepackage{fancyhdr}
\pagestyle{fancy}
\fancyhf{}
\fancyhead[L]{\small\color{litblue}Froth — Apuntes de ML}
\fancyhead[R]{\small\thepage}
\renewcommand{\headrulewidth}{0.4pt}
\usepackage[colorlinks=true,linkcolor=litblue,urlcolor=litblue]{hyperref}

% Cabecera de cada apunte: \cabecera{numero}{titulo}
\newcommand{\cabecera}[2]{%
  \begin{tcolorbox}[colback=litblue,coltext=white,colframe=litblue,arc=2mm]
    {\Large\bfseries Apunte #1 — #2}\\[3pt]
    {\small Proyecto Froth · aprender Machine Learning construyendo}
  \end{tcolorbox}\vspace{0.8em}}

% Cajas temáticas
\newtcolorbox{concepto}[1]{colback=blue!4,colframe=litblue,title={Concepto: #1},arc=1.5mm}
\newtcolorbox{practica}{colback=green!5,colframe=litgreen,title={Pruébalo tú (teclado en mano)},arc=1.5mm}
\newtcolorbox{ojo}{colback=orange!8,colframe=litorange,title={Ojo / error típico},arc=1.5mm}
\newtcolorbox{resumen}{colback=gray!8,colframe=gray!60,title={En una frase},arc=1.5mm}

\newcommand{\codigo}[1]{\texttt{#1}}
 justo después de \documentclass.
\usepackage[margin=2.4cm]{geometry}
\usepackage{xcolor}
\definecolor{litblue}{RGB}{31,79,121}
\definecolor{litgreen}{RGB}{27,94,32}
\definecolor{litorange}{RGB}{230,81,0}
\usepackage{parskip}
\usepackage{titlesec}
\titleformat{\section}{\Large\bfseries\color{litblue}}{\thesection}{0.6em}{}
\titleformat{\subsection}{\large\bfseries\color{litblue}}{\thesubsection}{0.6em}{}
\usepackage{tcolorbox}
\usepackage{fancyhdr}
\pagestyle{fancy}
\fancyhf{}
\fancyhead[L]{\small\color{litblue}Froth — Apuntes de ML}
\fancyhead[R]{\small\thepage}
\renewcommand{\headrulewidth}{0.4pt}
\usepackage[colorlinks=true,linkcolor=litblue,urlcolor=litblue]{hyperref}

% Cabecera de cada apunte: \cabecera{numero}{titulo}
\newcommand{\cabecera}[2]{%
  \begin{tcolorbox}[colback=litblue,coltext=white,colframe=litblue,arc=2mm]
    {\Large\bfseries Apunte #1 — #2}\\[3pt]
    {\small Proyecto Froth · aprender Machine Learning construyendo}
  \end{tcolorbox}\vspace{0.8em}}

% Cajas temáticas
\newtcolorbox{concepto}[1]{colback=blue!4,colframe=litblue,title={Concepto: #1},arc=1.5mm}
\newtcolorbox{practica}{colback=green!5,colframe=litgreen,title={Pruébalo tú (teclado en mano)},arc=1.5mm}
\newtcolorbox{ojo}{colback=orange!8,colframe=litorange,title={Ojo / error típico},arc=1.5mm}
\newtcolorbox{resumen}{colback=gray!8,colframe=gray!60,title={En una frase},arc=1.5mm}

\newcommand{\codigo}[1]{\texttt{#1}}

\begin{document}
\cabecera{37}{MMR: relevante pero distinto (y un objetivo que salió degenerado)}

\section{El problema que quedaba}

El centroide (apunte 35) tiene un defecto estructural: las frases que se parecen al
centro se parecen ENTRE SÍ. En tu cluster de geología, las 3 mejores frases tenían
similitud mutua de \textbf{0.945} --- un resumen diciendo lo mismo tres veces.

\begin{concepto}{Maximal Marginal Relevance (Carbonell \& Goldstein, 1998)}
Selección voraz, una frase a la vez. Cada candidata se puntúa como
\[
\lambda \cdot \mathrm{relevancia} \;-\; (1-\lambda) \cdot \mathrm{redundancia}
\]
donde redundancia = su similitud máxima con lo YA elegido. Con $\lambda = 1$ vuelves al
centroide puro; bajando $\lambda$ compras diversidad pagando relevancia. La primera
elegida siempre es la más relevante (no hay nada elegido aún con qué ser redundante).
\end{concepto}

\section{El tropiezo (que enseña más que el acierto)}

Primer intento de elegir $\lambda$ con datos: ``maximizar relevancia $-$ redundancia''.
Resultado: recomendó $\lambda = 0.5$... el BORDE de la grilla. Ese objetivo es
\textbf{degenerado}: bajar $\lambda$ SIEMPRE reduce la redundancia más rápido de lo que
pierde relevancia, así que el ``óptimo'' se escapa hacia abajo indefinidamente. Un
recomendador que siempre elige el extremo no está midiendo nada.

\section{La regla que sí funciona}

Mirando tu barrida real (corpus Beauvoir, 12 clusters):

\begin{center}
\begin{tabular}{cccc}
$\lambda$ & relevancia & redundancia \\
\hline
0.5 & 0.868 & 0.758 \\
0.6 & 0.906 & 0.829 \\
\textbf{0.7} & \textbf{0.918} & \textbf{0.860} \\
0.8 & 0.921 & 0.873 \\
1.0 & 0.923 & 0.884 \\
\end{tabular}
\end{center}

De 1.0 a 0.7 la relevancia casi no cae ($-0.5\%$) y la diversidad llega gratis; de 0.7
hacia abajo empiezas a pagar de verdad. Regla adoptada (interpretable): \textbf{el
$\lambda$ más diverso que conserve $\geq 99\%$ de la relevancia alcanzable}. Tu corpus
votó $\lambda = 0.7$ solo --- y coincide con el default de la literatura, pero ahora
SABEMOS por qué vale para estos datos (regla 9 del proyecto).

\begin{practica}
\texttt{.venv\textbackslash Scripts\textbackslash python.exe -m froth.review} ---
al final imprime la barrida y el antes/después del cluster más redundante: el centroide
da 3 frases clónicas de pegmatitas; MMR mantiene la definición y añade granitos Nb-Ta y
el Erzgebirge. Tres ángulos, no un eco.
\end{practica}

\begin{ojo}
MMR es voraz: elige lo mejor AHORA sin mirar adelante, porque el óptimo exacto es
combinatorio (probar todas las combinaciones de frases). Para resúmenes de 3-5 frases,
la aproximación voraz es el estándar y sobra.
\end{ojo}

\begin{resumen}
MMR = $\lambda \cdot$relevancia $-$ $(1-\lambda) \cdot$redundancia, selección voraz.
El primer objetivo para elegir $\lambda$ salió degenerado (se clavó al borde); la regla
final --- máxima diversidad conservando el 99\% de la relevancia --- eligió 0.7 con tus
datos. Con esto están las 3 piezas del ranking: típico (centroide), respaldado
(TextRank) y diverso (MMR). Sigue ensamblar el draft v2 (7.4).
\end{resumen}

\end{document}
